\documentclass[11pt,a4paper]{article}

% ── Packages ──
\usepackage[utf8]{inputenc}
\usepackage[T1]{fontenc}
\usepackage{amsmath,amssymb}
\usepackage{siunitx}
\usepackage{booktabs}
\usepackage{geometry}
\usepackage{hyperref}
\usepackage{enumitem}
\usepackage{float}

\geometry{margin=2.5cm}
\hypersetup{colorlinks=true,linkcolor=blue,citecolor=blue,urlcolor=blue}

\title{MARS Actuation System:\\
       Dual-Concentric Magnetic Drive with Differential Belt}
\author{}
\date{}

\begin{document}
\maketitle

% ══════════════════════════════════════════════════════════════════════════════
\section{Design Requirements}
\label{sec:requirements}
% ══════════════════════════════════════════════════════════════════════════════

The MARS propeller arm (Geometry~A) executes two sequential motions per
event:
\begin{enumerate}[nosep]
  \item \textbf{Phase~1 --- Azimuthal rotation ($\theta$):}
        The arm sweeps through $\theta_\mathrm{max} = \pi/N$ in a
        bang-bang profile. Required torque: $\tau = \SI{60}{N\cdot m}$.
        Peak tip velocity: $V_\mathrm{tip} \approx \SI{6}{m/s}$
        ($N = 8$). Duration: $t_\mathrm{rot} \approx \SI{400}{ms}$.
  \item \textbf{Phase~2 --- Radial translation ($r$):}
        The carrier plate slides along the arm from its parked position
        to the target radius. Required force: $F = \SI{20}{N}$.
        Travel: $\Delta r \le \SI{0.8}{m}$. Peak velocity:
        $V_\mathrm{max} \approx \SI{8}{m/s}$.
        Duration: $t_\mathrm{slide} \approx \SI{200}{ms}$.
\end{enumerate}

The actuation mechanism must satisfy three constraints:
\begin{itemize}[nosep]
  \item All motors, bearings, and lubricated components must reside
        \emph{outside} the xenon pressure boundary to preserve
        radiopurity of the active volume.
  \item No rotating or sliding seals may penetrate the pressure
        boundary (leak risk at \SI{15}{bar}).
  \item The two motions must be independently controllable and
        \emph{sequentially decoupled}: the carrier must not move during
        arm rotation, and vice versa.
\end{itemize}


% ══════════════════════════════════════════════════════════════════════════════
\section{Architecture Overview}
\label{sec:architecture}
% ══════════════════════════════════════════════════════════════════════════════

The design uses a \textbf{dual-concentric magnetic coupling} through a
static titanium re-entrant well, combined with a
\textbf{differential co-rotation} scheme that decouples the two degrees
of freedom without any clutch or brake.

The key components, from bottom to top:
\begin{enumerate}[nosep]
  \item A titanium re-entrant well protruding below the vessel bottom
        flange.
  \item Two concentric magnetic couplings (outer and inner) that
        transmit torque through the static Ti wall.
  \item Two concentric shafts rising through the copper floor shield
        into the xenon volume.
  \item A bevel-gear redirector at the arm hub that converts inner-shaft
        rotation into belt motion along the arm.
  \item A timing belt running the length of the arm, carrying the
        carrier plate.
\end{enumerate}


% ══════════════════════════════════════════════════════════════════════════════
\section{Re-Entrant Well and Magnetic Drive}
\label{sec:well}
% ══════════════════════════════════════════════════════════════════════════════

\subsection{Re-Entrant Well}

A cylindrical titanium (Ti-6Al-4V) well is welded to the vessel bottom
flange and protrudes \SI{200}{mm} downward into the atmosphere. The well
is a static part of the pressure boundary --- it has no moving parts, no
seals, and no feedthroughs. Its sole function is to bring the magnetic
coupling close enough for effective torque transmission.

\begin{table}[H]
\centering
\caption{Re-entrant well parameters.}
\label{tab:well}
\begin{tabular}{lll}
\toprule
Parameter & Value & Note \\
\midrule
Material             & Ti-6Al-4V & Non-ferromagnetic, low radioactivity \\
Outer diameter       & \SI{120}{mm} & Houses the \SI{100}{mm} outer coupling \\
Wall thickness       & \SI{4}{mm} & Pressure boundary \\
Length               & \SI{200}{mm} & Below bottom flange \\
Net internal pressure & \SI{14}{bar} & Xe at 15~bar minus atmosphere \\
Hoop stress          & \SI{21}{MPa} & $\sigma_h = P\,r/t$ \\
Safety factor        & 42 & vs.\ Ti-6Al-4V yield (\SI{880}{MPa}) \\
\bottomrule
\end{tabular}
\end{table}

Titanium is chosen because it is non-ferromagnetic (the magnetic flux
from the NdFeB magnets passes through the wall with minimal attenuation),
structurally strong, and has low intrinsic radioactivity.

\subsection{Magnetic Couplings}

Two independent synchronous magnetic couplings transmit torque through
the Ti wall without mechanical contact:

\paragraph{Outer coupling ($\theta$ drive).}
A ring of NdFeB permanent magnets on the external stator (air side)
locks magnetically to a corresponding ring on the internal rotor
(xenon side). The coupling must transmit up to $\tau = \SI{60}{N\cdot m}$
for the arm rotation. At a coupling diameter of \SI{100}{mm} and a total
radial air gap of $\sim\SI{6}{mm}$ (Ti wall + mechanical clearances),
this is within the rated capacity of commercial rare-earth magnetic
couplings.

\paragraph{Inner coupling ($r$ drive).}
A smaller set of NdFeB magnets, nested concentrically inside the outer
coupling, transmits torque for the radial belt drive. The required
torque is only $\tau_\mathrm{inner} = \SI{0.5}{N\cdot m}$
(Section~\ref{sec:belt}), which is trivial for any NdFeB coupling
at \SI{40}{mm} diameter.

Each coupling is driven by its own servo motor mounted on the air side
of the well.


% ══════════════════════════════════════════════════════════════════════════════
\section{Dual-Shaft Arrangement}
\label{sec:shafts}
% ══════════════════════════════════════════════════════════════════════════════

The two magnetic rotors drive two concentric shafts that rise vertically
through a bore in the copper floor shield:

\paragraph{Outer shaft ($\theta$).}
A hollow Ti tube ($\varnothing\,\SI{40}{mm}$ OD, \SI{3}{mm} wall)
connected to the outer magnetic rotor. It terminates above the floor
shield in a central \textbf{yoke} to which the NACA\,0012 arms are
bolted. Rotation of this shaft directly sets the azimuthal coordinate
$\theta$.

\paragraph{Inner shaft ($r$).}
A solid Ti rod ($\varnothing\,\SI{15}{mm}$) running concentrically
inside the outer shaft on PEEK sleeve bearings. It is connected to the
inner magnetic rotor and extends above the yoke, where it carries a
bevel gear (Section~\ref{sec:belt}).

\begin{table}[H]
\centering
\caption{Shaft parameters.}
\label{tab:shafts}
\begin{tabular}{lcc}
\toprule
& Outer shaft & Inner shaft \\
\midrule
Function           & Arm rotation ($\theta$) & Carrier translation ($r$) \\
Material           & Ti-6Al-4V & Ti-6Al-4V \\
Outer diameter     & \SI{40}{mm} & \SI{15}{mm} \\
Wall thickness     & \SI{3}{mm} (hollow) & solid \\
Max torque         & \SI{60}{N\cdot m} & \SI{0.5}{N\cdot m} \\
Shear stress       & \SI{10}{MPa} & $<\SI{1}{MPa}$ \\
Safety factor      & $>50$ & $>500$ \\
Bearings           & --- & PEEK sleeves \\
\bottomrule
\end{tabular}
\end{table}

The copper floor shield bore is $\varnothing\,\SI{120}{mm}$ to
accommodate the outer coupling housing. PEEK is used for the
inter-shaft bearings because of its low outgassing, radiation
tolerance, and good tribological properties in inert-gas environments.


% ══════════════════════════════════════════════════════════════════════════════
\section{Differential Co-Rotation Principle}
\label{sec:differential}
% ══════════════════════════════════════════════════════════════════════════════

The two degrees of freedom ($\theta$ and $r$) are decoupled by
exploiting the \emph{differential} angular velocity between the two
shafts. Define:
\begin{equation}
  \Delta\omega \;\equiv\; \omega_\mathrm{inner} - \omega_\mathrm{outer}\,.
  \label{eq:domega}
\end{equation}

The belt drive along the arm (Section~\ref{sec:belt}) is actuated by
$\Delta\omega$: only the \emph{relative} rotation between inner and
outer shafts moves the carrier. This yields a clean separation of the
two phases:

\paragraph{Phase~1 (arm rotation).}
Both servo motors drive their magnetic couplings \emph{synchronously}
at the same angular velocity profile:
\begin{equation}
  \omega_\mathrm{inner}(t) = \omega_\mathrm{outer}(t)
  \quad\Longrightarrow\quad
  \Delta\omega = 0\,.
\end{equation}
The bevel gear sees no input rotation, the belt is stationary relative
to the arm, and the carrier does not move. The entire assembly ---
both shafts, yoke, arms, belt, and carrier --- rotates as a rigid body.

\paragraph{Phase~2 (carrier translation).}
The outer motor holds the outer shaft at a fixed angle (arm stationary):
\begin{equation}
  \omega_\mathrm{outer} = 0\,,\quad
  \omega_\mathrm{inner} \ne 0
  \quad\Longrightarrow\quad
  \Delta\omega = \omega_\mathrm{inner}\,.
\end{equation}
Now the inner shaft rotates inside the stationary outer shaft. The
bevel gear converts this differential rotation into belt motion, and the
carrier translates radially along the arm.

This scheme requires no clutch, brake, or switchable mechanism inside
the xenon volume. The decoupling is purely kinematic, enforced by the
motor controllers.


% ══════════════════════════════════════════════════════════════════════════════
\section{Rotary-to-Linear Conversion: Bevel Gear and Timing Belt}
\label{sec:belt}
% ══════════════════════════════════════════════════════════════════════════════

\subsection{Bevel Gear Redirector}

The inner shaft extends above the yoke and terminates in a straight
bevel gear that redirects the rotation axis by $90^\circ$: from vertical
(shaft axis) to horizontal (along the arm). A mating bevel gear on a
short horizontal stub shaft drives the timing-belt pulley at the arm
root. A 1:1 gear ratio is used for simplicity.

\subsection{Timing Belt Along the Arm}

A toothed timing belt (e.g.\ GT3 profile, \SI{15}{mm} wide) runs the
length of each arm:
\begin{itemize}[nosep]
  \item \textbf{Drive pulley} ($r_p = \SI{25}{mm}$): at the arm root,
        driven by the bevel gear.
  \item \textbf{Idler pulley}: at or near the arm tip.
  \item \textbf{Belt path}: the belt forms a closed loop along two
        parallel channels inside the arm structure. Total belt length
        $\approx 2R + \pi(2 r_p) \approx \SI{3.4}{m}$.
  \item \textbf{Carrier attachment}: the carrier plate is clamped to
        one run of the belt. When the belt moves, the carrier translates
        radially.
\end{itemize}

\subsection{Kinematics}

The carrier linear velocity is set by the belt speed:
\begin{equation}
  V_\mathrm{carrier} = \Delta\omega\; r_p\,.
  \label{eq:Vcarrier}
\end{equation}

At peak carrier velocity $V_\mathrm{max} = \SI{8}{m/s}$:
\begin{align}
  \Delta\omega_\mathrm{max}
    &= \frac{V_\mathrm{max}}{r_p}
     = \frac{8.0}{0.025}
     = \SI{320}{rad/s}
     \approx \SI{3060}{rpm}\,.
  \label{eq:omega_inner}
\end{align}
This is well within the speed range of standard servo motors
(typical: 3000--6000~rpm rated).

\subsection{Force and Torque}

The carrier force is transmitted through the belt:
\begin{equation}
  F_\mathrm{carrier} = \frac{\tau_\mathrm{inner}}{r_p}\,,
  \label{eq:Fcarrier}
\end{equation}
so the required inner-shaft torque is:
\begin{equation}
  \tau_\mathrm{inner}
    = F_\mathrm{carrier} \times r_p
    = \SI{20}{N} \times \SI{0.025}{m}
    = \SI{0.5}{N\cdot m}\,.
  \label{eq:tau_inner}
\end{equation}

The belt tensions are modest:
\begin{equation}
  T_\mathrm{tight} - T_\mathrm{slack} = F_\mathrm{carrier}
    = \SI{20}{N}\,,
\end{equation}
with typical tight-side tension $T_\mathrm{tight} \approx \SI{24}{N}$
--- negligible for any standard timing belt.


% ══════════════════════════════════════════════════════════════════════════════
\section{Summary of Parameters}
\label{sec:summary}
% ══════════════════════════════════════════════════════════════════════════════

\begin{table}[H]
\centering
\caption{Actuation system --- key parameters.}
\label{tab:summary}
\begin{tabular}{llcc}
\toprule
& & Phase~1 ($\theta$) & Phase~2 ($r$) \\
\midrule
\multicolumn{4}{l}{\textit{Motor requirements}} \\
Driving shaft           & & Outer & Inner \\
Coupling torque         & & \SI{60}{N\cdot m} & \SI{0.5}{N\cdot m} \\
Coupling diameter       & & \SI{100}{mm} & \SI{40}{mm} \\
Peak shaft speed        & & $\sim$\SI{16}{rad/s} & \SI{320}{rad/s} \\
                        & & ($\sim$\SI{150}{rpm}) & ($\sim$\SI{3060}{rpm}) \\
\midrule
\multicolumn{4}{l}{\textit{Differential condition}} \\
$\omega_\mathrm{outer}$ & & $\omega(t)$ & 0 \\
$\omega_\mathrm{inner}$ & & $\omega(t)$ & $\Delta\omega(t)$ \\
$\Delta\omega$          & & 0 (belt locked) & $\ne 0$ (belt moves) \\
\midrule
\multicolumn{4}{l}{\textit{Belt drive (Phase~2 only)}} \\
Drive pulley radius     & $r_p$ & --- & \SI{25}{mm} \\
Bevel gear ratio        & & --- & 1:1 \\
Peak belt speed         & & --- & \SI{8}{m/s} \\
Belt tension (tight)    & & --- & \SI{24}{N} \\
Carrier force           & & --- & \SI{20}{N} \\
Carrier travel          & & --- & $\le\SI{0.8}{m}$ \\
Slide time              & & --- & $\sim$\SI{200}{ms} \\
\midrule
\multicolumn{4}{l}{\textit{Structural margins}} \\
Ti well hoop stress / yield & & \multicolumn{2}{c}{\SI{21}{MPa}\,/\,\SI{880}{MPa} (SF\,$=$\,42)} \\
Outer shaft shear / yield   & & \multicolumn{2}{c}{\SI{10}{MPa}\,/\,\SI{550}{MPa} (SF\,$>$\,50)} \\
\bottomrule
\end{tabular}
\end{table}


% ══════════════════════════════════════════════════════════════════════════════
\section{Discussion}
\label{sec:discussion}
% ══════════════════════════════════════════════════════════════════════════════

\paragraph{Why outer = $\theta$, inner = $r$.}
The arm rotation requires the high torque (\SI{60}{N\cdot m}), while
the belt drive requires only \SI{0.5}{N\cdot m}. The outer magnetic
coupling has the larger diameter and hence higher torque capacity per
unit length, making it the natural choice for the high-torque channel.
The inner coupling, though smaller, easily handles the belt-drive torque.
Swapping the roles would require the outer shaft to extend above the arm
attachment point with an additional bearing stage, adding mechanical
complexity for no benefit.

\paragraph{Servo control.}
The co-rotation requirement during Phase~1 demands that both motors
track the same angular profile to within a small tolerance. If the
inner motor lags or leads by $\Delta\theta_\mathrm{err}$, the carrier
shifts by $\Delta r = r_p\,\Delta\theta_\mathrm{err}$. For a
\SI{25}{mm} pulley, a $1^\circ$ tracking error produces
$\Delta r \approx \SI{0.4}{mm}$ --- acceptable, and well within
the capability of modern servo drives with encoder feedback.

\paragraph{Radiopurity.}
All motors, magnets, bearings, and lubricants are outside the pressure
boundary or confined to the well interior (PEEK bearings in the shaft
annulus). No organic lubricants or outgassing materials are present
in the active xenon volume above the copper floor shield. The NdFeB
magnets inside the well (xenon side) can be encapsulated in
electroless-nickel or Ti cladding to prevent direct contact with xenon.

\paragraph{Belt in the arm structure.}
The timing belt runs inside channels in the arm (which has an
aerodynamic NACA\,0012 external profile but can be hollow or
truss-structured internally). The belt adds negligible mass
($\sim$\SI{50}{g} for a \SI{3.4}{m} GT3 belt) and does not affect
the external flow geometry.

\end{document}
